\uuid{WBCr}
\exo7id{1836}
\titre{exo7 1836}
\auteur{gourio}
\organisation{exo7}
\datecreate{2001-09-01}
\isIndication{false}
\isCorrection{false}
\chapitre{Extremum, extremum lié}
\sousChapitre{Extremum, extremum lié}
\cpge{mpsi-mp2i.convexite.b, mpsi-mp2i.fonctions-deux-variables.b, pcsi.fonctions-deux-variables.b, ptsi.fonctions-deux-variables.b, mpsi-mp2i.fonctions-deux-variables.d, pcsi.fonctions-deux-variables.d, ptsi.fonctions-deux-variables.d, bcpst1.derivees-primitives-integrales.b, bcpst1.fonctions-deux-variables.c, mp-mpi.calcul-differentiel-optimisation.f}
\module{Calcul différentiel}
\niveau{L3}
\difficulte{}

\contenu{
\question{
Si $ f$ est concave sur un ouvert convexe $ U\subset {\Rr}^{2}$ et si :
$$\exists  a\in U,\frac{\partial f}{\partial x_{1}}(a)=\frac{\partial f}{%
\partial x_{2}}(a)=0, $$
alors $ f$ admet un maximum local en $ a.$
}
}
